\documentclass[10pt,a4paper]{amsart}
\usepackage[margin=19mm]{geometry}
\usepackage{amsmath,amssymb,mathtools}
\usepackage[scr=rsfs,cal=euler]{mathalfa}
\usepackage{xfrac}
\usepackage[unicode,hidelinks,bookmarks=false]{hyperref}
\newcommand{\ie}{i.e.\ }
\newcommand{\cf}{cf.\ }
\newcommand{\resp}{respectively\ }
\newcommand{\Subsection}{\subsection}
\providecommand{\nobreakdash}{\nobreak\hbox{-}}
\setlength{\emergencystretch}{2em}
\multlinegap0pt
\newtheorem{comment}{Comment}
\usepackage{dsfont}
\usepackage{bbm}
\usepackage{enumerate}
\usepackage{amssymb}
\usepackage{xcolor}
\usepackage{tikz}
\usepackage{mathtools}
\usepackage{float}
\usepackage{bm}
\newtheorem{lem}{Lemma}
\newtheorem{prop}{Proposition}
\newcommand{\m}{\mathrm{m}}
\newcommand{\mb}{\mathbb}
\newcommand{\mf}{\mathfrak}
\newcommand{\n}{\mf{n}_D}
\newcommand{\ol}{\overline}
\title{\boldmath An analogue of the Dedekind Eta function \\ for Hecke Groups $H(\sqrt{D})$}
\author{Debmalya Basak, Dorian Goldfeld, Winston Heap, Nicolas Robles, Alexandru Zaharescu}
\date{Source: arXiv:2508.21239v3 (CC BY 4.0)}
\begin{document}
\maketitle

\noindent\emph{Source and licence.} \boldmath An analogue of the Dedekind Eta function \\ for Hecke Groups $H(\sqrt{D})$. Version arXiv:2508.21239v3 (CC BY 4.0), distributed by arXiv under the Creative Commons Attribution 4.0 International licence, http://creativecommons.org/licenses/by/4.0/, as recorded in the arXiv licence metadata for this exact version and shown on the abstract page https://arxiv.org/abs/2508.21239v3. Full licence notice: \texttt{licenses/arxiv-2508.21239-CC-BY-4.0.txt}.\par\smallskip

\noindent\textit{Editorial scope.} Counted unit: the article's prop lem (source0.tex lines 917--949). The article's complete original proof of it is delivered with it (source0.tex lines 952--1219, one \texttt{proof} environment of 268 lines). No mathematics is written, repaired or re-derived: the statement and the proof are the article's own lines, in the article's own notation, and no retained line is substituted or re-flowed.  Retained uncounted as the source-local context the proof needs: lines 568--570; lines 843--853.  Nothing was omitted from any retained span, so the package carries no omission record.  No retained line contains a citation, so the wrapper carries no bibliography and no bibliography entry.  No retained line contains a figure environment, an image-inclusion command or drawing code, so no image is delivered and the package declares no asset dependency.  Editorial formatting only, in addition: an AMS \texttt{amsart} preamble that restates the article's own declarations and macros used by the retained lines, namely the environments \texttt{lem}, \texttt{prop} on separate counters, together with the standard packages those lines need.  The retained environments print in this document as Lemma 1, Proposition 1 in order of appearance, whereas the article prints the same items with the numbering of its own full document.  The complete native source of the article as distributed in the arXiv e-print for this exact version is bundled unchanged as \texttt{sources/184095/source0.tex}.
\begin{align} \label{FunctionalEquationForL(s,chi)}
\Lambda(s,{\chi_{D}}) := \left(\frac{D}{\pi} \right)^{\tfrac{s}{2}} \Gamma\left(\frac{s}{2} \right) L(s,{\chi_{D}}) =  \Lambda(1-s,{\chi_{D}}).
\end{align}

\begin{lem}\label{lem 1} For $|q|<1$, we have 
\begin{equation}\label{log F}
\log G(q)
=
-\sum_{m\geq 1}\frac{g(q^m)}{m(1-q^{Dm})}
-\sqrt{D}\sum_{m\geq 1}\frac{\chi_D(m)}{m}\frac{q^m}{1-q^m}
\end{equation}
where $
g(x)=\sum_{a=1}^{D-1}\chi_D(a)x^a$
is the Fekete polynomial corresponding to the character $\chi_D$.
\end{lem}

\begin{prop}\label{precise lem}
Let $\varepsilon>0$. Fix $u,v \in \mathbb{N}$ satisfying $1 \leq u \leq v$ and $(u,v)=1$. Set $q=re^{2\pi i {u}/{v} +i\theta}$ where $0<r<1$ and $|\theta|\leq B(1-r)$ for some fixed constant $B>0$. If $D\nmid v$, we have 
\begin{equation}\label{log F 2}
\log G(q)
=
%\frac{\sqrt{5}L(2,\chi_D)}{4(\log(1/r)-i\theta)}
-\frac{\chi_D(v)\sqrt{D}L(2,\chi_D)}{v^2(\log(1/r)-i\theta)}
+\gamma_D(u,v)
+\mathcal{O}_{D,B,\varepsilon}(v^{1+\varepsilon}(1-r))
\end{equation}
where 
\begin{align}\label{Eq: Gamma Def}
\gamma_D(u,v)
&=
\chi_D(v) L(1,\chi_D) \frac{(1-v)\sqrt{D}}{2v}
\notag \\
&\quad -
\frac{1}{\phi(v)}\sum_{\psi \bmod v} \sum_{a=1}^{v-1}\overline{\psi}(a)\Big[\sqrt{D}P(0,au/v)L(1,\chi_D\psi)
+
P(1,au/v)L(0,\chi_D\psi)\Big].
\end{align}
%and $P(s,\alpha)$ is the  periodic zeta function. 

If $D|v$, we have 
\begin{equation}\label{log F 3}
\log G(q)
=
%\frac{\sqrt{5}L(2,\chi_D)}{4(\log(1/r)-i\theta)}
-\frac{\chi_D(u)D^{3/2}L(2,\chi_D)}{v^2(\log(1/r)-i\theta)}
+a_D(u,v)\log(\log(1/r)-i\theta)+b_D(u,v)+\mathcal{O}_{D,B,\varepsilon}(v^{1+\varepsilon}(1-r))
\end{equation}
for some computable constants $a_D(u,v),b_D(u,v)$ depending only on $u,v$ and $D$. 
\end{prop}

\begin{proof}
For notational simplicity, let us replace $q$ by $qe(u/v)$ where we set $q=re^{i\theta}$ and $|\theta|\leq B(1-r)$. Recall from the proof of Lemma \ref{lem 1} that we may write
\begin{equation}\label{log F twist}
\log G(qe(u/v))=-S_1-\sqrt{D}S_2
\end{equation}
with 
\[
S_1
=
\sum_{m,n\geq 1}\frac{\chi_D(n)e(mnu/v)q^{mn}}{m}
\qquad \textrm{and} \qquad
S_2
=
\sum_{m,n\geq 1}\frac{\chi_D(m)e(mnu/v)q^{mn}}{m}.
\]
We break our proof into the following parts.\\

 \noindent \textbf{$S_2$ Analysis.} We look at $S_2$ first and consider the sum over $m$ modulo $v$. Write $S_2=S_{21}+S_{22}$ with 
 \[
 S_{21}
 =
 \frac{\chi_D(v)}{v}\sum_{m,n\geq 1}\frac{\chi_D(m)}{m}q^{vmn},
 \qquad \textrm{and} \qquad
 S_{22}
 =
 \sum_{a=1}^{v-1}\sum_{\substack{m,n\geq 1 \\ m\equiv a \bmod v}}
 \frac{\chi_D(m)e(anu/v)q^{mn}}{m}.
 \]

Applying the identity
\begin{equation}\label{mellin}
\qquad e^{-z}=\frac{1}{2\pi i }\int_{c-i\infty}^{c+i\infty}\Gamma(s)z^{-s}\, ds, \qquad c>0,
\end{equation}
which is valid for $\operatorname{Re}(z)>0$, we have
\begin{align*}
S_{21} &=
\frac{\chi_D(v)}{v}
\sum_{m,n\geq 1}\frac{\chi_D(m)}{m}
\frac{1}{2\pi i }\int_{2-i\infty}^{2+i\infty}\Gamma(s)(vmn\log(1/q))^{-s} \, ds
\\
&=
\frac{\chi_D(v)}{v}
\frac{1}{2\pi i }\int_{2-i\infty}^{2+i\infty}
\Gamma(s)\zeta(s)L(1+s,\chi_D)(v\log(1/q))^{-s}\, ds.
\end{align*}
The integrand has simple poles at $s=1,0,-1$ and so we would like to shift the integral past these points. 

When estimating the integral in the left half-plane note that since $|\theta|\leq B(1-r)$, 
\[
|\arg \log(1/q)|=\arctan(|\theta|/\log(1/r))\leq \tfrac{\pi}{2}-\delta
\] 
for some fixed $\delta=\delta(B)\asymp \min (1,B^{-1})$. In particular, by Stirling's formula we have 
\[
\Gamma(s)(\log(1/q))^{-s}\ll |\log(1/q)|^{-\operatorname{Re}(s)} e^{-\delta \operatorname{Im}(s)/2}, \quad \textrm{as} \quad \operatorname{Im}(s)\to\infty.
\]
Hence the integral over the new line is absolutely convergent and $\ll_{D,B} |\log(1/q)|^{-\operatorname{Re}(s)}$ since $\zeta(s)$ and $L(1+s,\chi_D)$  are of polynomial growth at most.
Thus shifting to $\operatorname{Re}(s)=-3/2$, say, we find that 
\begin{align*}
S_{21}
= &
\frac{\chi_D(v)L(2,\chi_D)}{v^2\log(1/q)}
-\frac{\chi_D(v)}{2v}
L(1,\chi_D)
+\mathcal{O}_{D,B}(v^{1/2}|1-q|)
\end{align*}
on recalling that $\zeta(0)=-1/2$. 


By the orthogonality of Dirichlet characters, 
\[
S_{22}
=
\frac{1}{\phi(v)}\sum_{\psi \bmod v} \sum_{a=1}^{v-1}\overline{\psi}(a)\sum_{\substack{m,n\geq 1}}
 \frac{\chi_D(m)\psi(m)e(anu/v)q^{mn}}{m}.
\]
On applying \eqref{mellin}, this is equal to
\begin{align}\label{Eq: Double Pole}
\frac{1}{\phi(v)}\sum_{\psi \bmod v} \sum_{a=1}^{v-1}\overline{\psi}(a)
\frac{1}{2\pi i }\int_{2-i\infty}^{2+i\infty}
\Gamma(s)P(s,au/v)L(1+s,\chi_D\psi)(\log(1/q))^{-s}\, ds
\end{align}
where $P(s,\alpha)$ is the periodic zeta function. Since $au/v\in \mb{Q}\backslash\mb{Z}$, on shifting to the line $\operatorname{Re}(s)=-3/2$, we only encounter poles at $s=0,-1$. If $\chi_D\psi$ is principal, we have a double pole at $s=0$; otherwise it is simple. Note that if $\chi_D\psi$ is principal, we must have $D \mid v$. 
\begin{comment}
    Let us examine when this can occur. Writing $D=p_1\cdots p_r$ we see that $\chi_D=\chi_{p_1}\cdots \chi_{p_r}$ with $\chi_{p_i}$ the real character modulo $p_i$. Hence for $\chi_D\psi$ to be principal we must have $\psi=\chi_D\psi'$ with $\psi'$ principal and this can only occur if $D|v$. In this case  $\psi'$ is the principal character modulo $v/D$, or equivalently, $\psi=\chi_D\psi_0$ where $\psi_0$ is the prinicpal character modulo $v$, that is, $\psi$ is induced by $\chi_D$ and in this case $\chi_D\psi$ is the principal character modulo $v$. 
\end{comment}
Distinguishing between the two cases and accounting for the double pole when $D|v$, we have
\begin{align*}
S_{22}& =
\frac{\mathbbm{1}_{D\nmid v}}{\phi(v)}\sum_{\psi \bmod v} \sum_{a=1}^{v-1}\overline{\psi}(a)P(0,au/v)L(1,\chi_D\psi)
\\
&\quad +
\mathbbm{1}_{D| v}(a_D(u,v)\log\log(1/q)+b_D(u,v))
+
\mathcal{O}_{D,B}(v|1-q|),
\end{align*}
for some computable constants $a_D(u,v),b_D(u,v)$ depending only on $u,v$ and $D$. In total, we  find that 
\begin{align}\label{S2}
S_2
&=
\mathbbm{1}_{D\nmid v} \bigg(\frac{\chi_D(v)L(2,\chi_D)}{v^2\log(1/q)}
-\frac{\chi_D(v)}{2v}
L(1,\chi_D)
+
h_D(u,v)\bigg) \notag\\
&\quad +
\mathbbm{1}_{D| v}(a_D(u,v)\log\log(1/q)+b_D(u,v))
+\mathcal{O}_B(v|1-q|),
\end{align}
where 
\begin{align*}
h_D(u,v)
= &
\frac{1}{\phi(v)}\sum_{\psi \bmod v} \sum_{a=1}^{v-1}\overline{\psi}(a)P(0,au/v)L(1,\chi_D\psi).
%=
%\mathds{1}_{D\nmid v} \sum_{a=1}^{v-1}P(0,au/v)L(1,\chi_D, a)
\end{align*}
%and $L(s,\chi,a)=\sum_{n\equiv a\mod v}\chi(n)n^{-s}$.

\noindent \textbf{$S_1$ Analysis : The Case $D \nmid v$.} We consider the sum over $n$ modulo $v$ in $S_1$ and write $S_1=S_{11}+S_{12}$ with 
\[
S_{11}
=
\chi_D(v)\sum_{m,n\geq 1}\frac{\chi_D(n)q^{vmn}}{m}
,\qquad \textrm{and} \qquad
S_{12}
=
\sum_{a=1}^{v-1}\sum_{\substack{m,n\geq 1\\n\equiv a \bmod v}}\frac{\chi_D(n)e(amu/v)q^{mn}}{m}.
\] 
By \eqref{mellin}, we have
\[
S_{11}
=
\frac{\chi_D(v)}{2\pi i }\int_{2-i\infty}^{2+i\infty}
\Gamma(s)\zeta(1+s)L(s,\chi_D)(v\log(1/q))^{-s}\, ds.
\]
By the functional equation \eqref{FunctionalEquationForL(s,chi)}, $L(0,\chi_D)=0$ since $\chi_D$ is even. Thus the pole of the integrand at $s=0$ is simple. Shifting to $\operatorname{Re}(s)=-1-\varepsilon$ say, we see that
\begin{align}
 S_{11} & = \chi_D(v)L'(0,\chi_D)+\mathcal{O}_{B, \varepsilon}(v^{1+\varepsilon}|1-q|) \notag \\
 &=
  \frac{ \chi_D(v)\sqrt{D}}{2}L(1,\chi_D)+\mathcal{O}_{D,B, \varepsilon}(v^{1+\varepsilon}|1-q|), \label{S11}
 \end{align}
after another application of the functional equation. 



For $S_{12}$, by orthogonality of Dirichlet characters, we have 
\begin{align*}
S_{12}& =
\frac{1}{\phi(v)}\sum_{\psi \bmod v}\sum_{a=1}^{v-1}\ol{\psi}(a)
\sum_{\substack{m,n\geq 1}}\frac{\chi_D(n)\psi(n)e(amu/v)q^{mn}}{m}
\\
& =
\frac{1}{\phi(v)}\sum_{\psi \bmod v}\sum_{a=1}^{v-1} 
\frac{\ol{\psi}(a)}{2\pi i}\int_{2-i\infty}^{2+i\infty}
\Gamma(s)P(1+s,au/v)L(s,\chi_D\psi)(\log(1/q))^{-s}\, ds.
\end{align*}
As noted before, the character $\chi_D\psi$ can only be principal if $D|v$ which we are currently assuming is not the case. Hence we only have simple poles at $s=0,-1$ and on shifting contours,  
we see
\begin{equation}\label{S12}
S_{12}
=
%\mathds{1}_{D|v}\bigg[\frac{\prod_{p|v}(1-p^{-1})}{\phi(v)\log(1/q)}\sum_{a=1}^{v-1} \chi_D\psi_0(a)P(2,au/v)+c_{u,v}\bigg]
%\\
%+
\frac{\mathbbm{1}_{D\nmid v}}{\phi(v)}\sum_{\psi \bmod v}\sum_{a=1}^{v-1} 
{\ol{\psi}(a)}P(1,au/v)L(0,\chi_D\psi)
+ \mathcal{O}_{D,B}(v|1-q|).
\end{equation}
%Combining \eqref{S2},\eqref{S11} and \eqref{S12} in \eqref{log F twist} gives the result for $D\nmid v$. 



%Now, of course $\phi(v)^{-1}\prod_{p|v}(1-p^{-1})=v^{-1}$ whilst
%\[
%\sum_{a=1}^{v-1} \chi_D\psi_0(a)P(2,au/v)
%=
%\sum_{n\geq 1}\frac{1}{n^2}\sum_{a=1}^{v-1}\chi_D\psi_0(a)e(aun/v).
%\]
%Since $(u,v)=1$ the inner sum on the right is 
%\[
%\chi_D(u)\sum_{a=1}^{v-1}\chi_D\psi_0(a)e(an/v).
%\]
%If $(n,v)=1$ we may also factor out $\chi_D\psi_0(n)$ and remove the $n$ in the exponent of the summand. If $g:=(n,v)>1$ then $n/v=c/d$, say, with $(c,d)=1$ and 

\noindent \textbf{$S_1$ Analysis : The Case $D \mid v$.} Here we use a different decomposition of the sum. Let $v=Dv'$. We write
\[
S_1
=
\sum_{a=1}^{D-1}\chi_D(a)\sum_{\substack{m,n\geq 1\\n\equiv a \bmod D}}\frac{e(umn/v)q^{mn}}{m}
=
\sum_{a=1}^{D-1}\chi_D(a)\sum_{\substack{n\geq 0,m\geq 1}}\frac{e(aum/v)e(umn/v')q^{m(a+Dn)}}{m}.
\]
Now we consider the sum over $m$ modulo $v'$ giving 
\[
S_1=S_{11}'+S_{12}'
\]
with 
\[
S_{11}'
=
\frac{1}{v'}\sum_{a=1}^{D-1}\chi_D(a)\sum_{\substack{m\geq 1}}\frac{e(aum/D)}{m}\sum_{n\geq 0}q^{v'm(a+Dn)}
\]
and
\[
S_{12}'
=
\sum_{a=1}^{D-1}\chi_D(a)\sum_{b=1}^{v'-1}\sum_{\substack{m\geq 1\\m\equiv b \bmod v'}}\frac{e(aum/v)}{m}\sum_{n\geq 0}e(ubn/v')q^{m(a+Dn)}.
\]
Applying Mellin inversion and writing $v'(a+Dn)=v(n+a/D)$, we have 
\[
S_{11}'
=
\frac{1}{v'}\sum_{a=1}^{D-1}\chi_D(a)
\frac{1}{2\pi i }\int_{2-i\infty}^{2+i\infty}
\Gamma(s)P(1+s, au/D)\zeta(s,a/D)(v\log(1/q))^{-s}\, ds
\]
where $\zeta(s,\alpha)=\sum_{n\geq 0}(n+\alpha)^{-s}$ is the Hurwitz zeta function. This can be continued to a meromorphic function on $\mb{C}$ with only a simple pole at $s=1$ with residue $1$, and with at most polynomial growth in fixed vertical strips. Since $au/D$ is not integral, the periodic zeta function in the integrand is entire and so the poles at $s=1,0,-1$ are all simple. Shifting contours to $\operatorname{Re}(s)=-1-\varepsilon$ gives
\begin{align} \label{S11' Contribution}
S_{11}'
&=
\frac{D}{v^2\log(1/q)}\sum_{a=1}^D\chi_D(a)P(2,au/D) \notag \\
&\quad +
\sum_{a=1}^D\chi_D(a)P(1,au/D)\zeta(0,a/D)
+\mathcal{O}_{D,B, \varepsilon}(v^{\varepsilon}|1-q|)
\end{align} 
since $vv'=v^2/D$.
The constant in the leading term here, that is, the sum over $a$ involving $P(2,au/D)$ can be evaluated using Gauss sums as follows: 
\begin{align*}
\sum_{a=1}^D\chi_D(a)P(2,au/D) &=\sum_{n\geq 1}\frac{1}{n^2}\sum_{a=1}^D\chi_D(a)e(aun/D)\\
&=
\tau(\chi_D)\chi_D(u)\sum_{n\geq 1}\frac{\chi_D(n)}{n^2}
=
\sqrt{D}\chi_D(u)L(2,\chi_D).
\end{align*}



For $S_{12}'$, \eqref{mellin} gives 
\[
\sum_{a=1}^{D-1}\chi_D(a)\sum_{b=1}^{v'-1}
\frac{1}{2\pi i }\int_{2-i\infty}^{2+i\infty}
\Gamma(s)\bigg(\sum_{\substack{m\geq 1\\m\equiv b \bmod v'}}\frac{e(aum/v)}{m^{1+s}}\bigg)\bigg(\sum_{n\geq 0}\frac{e(ubn/v')}{(n+a/D)^s}\bigg)(D\log(1/q))^{-s}\, ds.
\]
The Dirichlet series in the integrand here are given by Lerch zeta functions. Since neither $au/v$ or $ub/v'$ are integral, these are entire functions and we only pick up poles at $s=0,-1$. Therefore, we obtain
\begin{align}
\label{S12' Contribution}
S_{12}' = c_D(u,v) +\mathcal{O}_{D,B, \varepsilon}(v^{\varepsilon} \lvert 1-q \rvert),
\end{align}
for some computable constant $c_D(u,v)$ depending only on $u,v$ and $D$.\\

\noindent \textbf{Conclusion.} We now put together the cases $D|v$ and $D\nmid v$ for $S_1$. Combining \eqref{S11}--\eqref{S12' Contribution}, we have
\begin{align}\label{S1 Final}
S_1
&=
  \mathbbm{1}_{D\nmid v} \bigg(\frac{ \chi_D(v)\sqrt{D}}{2}L(1,\chi_D)+h_D'(u,v)\bigg)
\notag \\
&\quad +
\mathbbm{1}_{D|v}\bigg(\frac{\chi_D(u)D^{3/2}L(2,\chi_D)}{v^2\log(1/q)}+c_D'(u,v)\bigg)+\mathcal{O}_{D,B, \varepsilon}(v^{1+\varepsilon}|1-q|)
\end{align}
for some computable constant $c_D'(u,v)$ where 
\[
h_D'(u,v)
=
\frac{1}{\phi(v)}\sum_{\psi \bmod v}\sum_{a=1}^{v-1} 
{\ol{\psi}(a)}P(1,au/v)L(0,\chi_D\psi).
\] 
For clarity, we mention that the second term on the right hand side of \eqref{S11' Contribution} has been absorbed into $c_D'(u,v)$. Finally, we note that in our range of $\theta$, $\lvert 1-q \rvert \asymp_B 1-r.$ Therefore, combining \eqref{S1 Final} with \eqref{S2} in \eqref{log F twist} and readjusting the constants $a_D(u,v),b_D(u,v)$ accordingly gives the desired result.
\end{proof}

\end{document}
